% Chapter 3, section 3: Request and reply
% Write in plain LaTeX. Code of a guide step: \guide{<step>}; commands: \begin{shell}...\end{shell};
% a screenshot: \screenshot{file}{caption}; a gap in the draft: \placeholder{...}.
% See book/latex/queen.sty for the full list.
%
% Draft by Claude, 2026-10-24. The code is cut from examples/tutorials/js/book-use-cases.mjs, which
% checks every result quoted here against a 2.6.0 node.

\section{Request and reply}

The booking page needs a price from the pricing service, now. If no answer comes in half a second,
it shows ``ask us''. A durable queue is too heavy for that, and a direct HTTP call ties the two
services together.

Ephemeral queues live in memory, outside the log. The caller makes a private inbox, sends it as
\texttt{replyTo}, and waits for the answer with a timeout. Any instance of the pricing service can
answer.

\queenfig{request-reply}{A quote with a deadline: the answer, or "ask us".}

\guide{uc-reply}

The answer came back in time: \texttt{\{room: 'deluxe', price: 129\}}. With no pricing service
running, the caller gave up after half a second and got \texttt{price: null}.

\why{an answer that comes too late is worth nothing, so it does not need a disk. In memory a round
trip costs very little, and the service scales like any consumer group.}

One thing to know: a node that crashes loses the ephemeral messages it held. There are no
transactions, dead letters or replays on these queues.
